% GENERATED FILE. Edit canonical lessons, then run npm run generate:books.
% canonical-source-hash: fnv1a64:5f42597a2da4f798
% canonical-lessons: FR-C20-je-suis-desole

\chapter{Sorry}
\label{ch:fr-sorry}
\hlchaptermodality{20}

\begin{chapteropening}
\textbf{By the end of this chapter:} \emph{I can apologise in French and pick between je suis désolé(e), pardon and excusez-moi for the situation I am actually in.}

Say je suis désolé(e) knowing it literally means 'I am desolate', then place pardon as the quick reflex and excusez-moi as the attention-getter. The chapter's only lesson, so the payoff covers everything it introduces.
\end{chapteropening}

\section[je suis désolé(e)]{je suis désolé(e) — \textquotedblleft{}I'm sorry\textquotedblright{} (literally \textquotedblleft{}I am desolate\textquotedblright{})}

\label{lesson:FR-C20-je-suis-desole}



\begin{quote}
You already know \textbf{s'il vous plaît}, \textquotedblleft{}please\textquotedblright{} — \textquotedblleft{}if it pleases you.\textquotedblright{} French's \textquotedblleft{}sorry\textquotedblright{} is, secretly, a much bigger word than it sounds: \textbf{désolé} is the same word as English \textbf{desolate}.
\end{quote}

\begin{cousinweb}
\begin{itemize}
\raggedright
  \item \textbf{je suis} = \textquotedblleft{}\textbf{I am}.\textquotedblright{}
  \item \textbf{désolé} (feminine \textbf{désolée}) = \textquotedblleft{}\textbf{sorry},\textquotedblright{} but literally \textquotedblleft{}\textbf{made utterly alone / devastated.}\textquotedblright{} It comes from Latin \textbf{dēsōlāre}, \textquotedblleft{}to make solitary\textquotedblright{} — \textbf{dē-} (\textquotedblleft{}thoroughly, completely\textquotedblright{}) + \textbf{sōlus} (\textquotedblleft{}alone\textquotedblright{}), the very root of English \textbf{sole}, \textbf{solo}, \textbf{solitary}, and \textbf{desolate} itself.
\end{itemize}

So \textbf{je suis désolé(e)} literally says \textquotedblleft{}\textbf{I am desolate}\textquotedblright{} — French reaches for a word of real devastation to say a plain \textquotedblleft{}sorry,\textquotedblright{} the way English keeps \textquotedblleft{}terribly sorry\textquotedblright{} as an option rather than a requirement.
\end{cousinweb}

\begin{culture}
French keeps two more courtesy words for different jobs:

\begin{itemize}
\raggedright
  \item \textbf{pardon} — a quick, all-purpose \textquotedblleft{}sorry / pardon me\textquotedblright{} (bumping someone, a small slip). It comes from \textbf{pardonner}, Latin \textbf{per-} (\textquotedblleft{}thoroughly\textquotedblright{}) + \textbf{dōnāre} (\textquotedblleft{}to give\textquotedblright{}) — \textquotedblleft{}to give completely,\textquotedblright{} i.e. to forgive. The neat part isn't that French and Spanish coined this separately — it's that they \textbf{didn't have to}: \textbf{perdōnāre} was already one word in Late Latin, so French \emph{pardonner} and Spanish \emph{perdón/perdonar} are simply two children of the \emph{same} inherited word, not two languages independently landing on the same idea.
  \item \textbf{excusez-moi} — \textquotedblleft{}\textbf{excuse me},\textquotedblright{} from \textbf{excuser} < Latin \textbf{ex-} (\textquotedblleft{}out of, free from\textquotedblright{}) + \textbf{causa} (\textquotedblleft{}cause, charge\textquotedblright{}) — literally \textquotedblleft{}to free from a charge.\textquotedblright{} Reach for this to get someone's \textbf{attention}, not to apologize for a real fault.
\end{itemize}

\textbf{Je suis désolé(e)} carries the most \textbf{regret}; \textbf{pardon} is the quick, light reflex; \textbf{excusez-moi} is for getting attention. Native speakers slide between all three by feel.
\end{culture}

\subsection*{Your turn}
Say these aloud:

\begin{itemize}
\raggedright
  \item \textquotedblleft{}désolé\textquotedblright{} — sorry / desolate — then \textquotedblleft{}je suis désolé,\textquotedblright{} I am sorry
  \item the quick one — \textquotedblleft{}pardon\textquotedblright{}
  \item the attention one — \textquotedblleft{}excusez-moi\textquotedblright{}
\end{itemize}

\begin{tcolorbox}[breakable,colback=teal!4,colframe=teal!35!black,title={Before you move on}]
What does \textbf{désolé} literally mean, and what English word is it identical to? (\textbf{\textquotedblleft{}Made utterly alone\textquotedblright{}} — the same word as English \textbf{desolate}.) What is \textbf{pardon} built from, and why does \textbf{Spanish perdón} look just like it? (\textbf{Per- \textquotedblleft{}thoroughly\textquotedblright{} + dōnāre \textquotedblleft{}to give\textquotedblright{}} — both simply inherited the \emph{same} already-formed Latin word, \textbf{perdōnāre}.) What's \textbf{excusez-moi} for? (\textbf{Getting someone's attention}, not apologizing for a real fault.)
\end{tcolorbox}
